% !TEX root = ../cubic-moduli-descent.tex
\section{点曲線の積における全中心同変核と混合型}
\label{mbfc:sec:full-comparison}

本節では、三ブロック以上、または二つの Jacobian が幾何学的に非同種である場合に、
第\ref{mbpd:sec:picard-descent}節の Picard 補正型を全ての中心同変比較核へ拡張する。
三ブロック以上では、指定した双対中心作用が因子を分離するため、同種性を除外しない。
二つの因子が同種の場合には、実際に外積でない混合核が生じる。
一般の体では、偏極曲線自己同型による追加の Kummer 補正を許す必要がある。
一方、固定 Hesse level を持つ族ではこの追加補正が消えるため、
系\ref{mbpd:cor:phase-independence}の最小次数は全ての中心同変比較核について成り立つ。

以下では $K\supset\C$ を体、$k\ge2$ とし、
第\paperref{main}{mbrc:sec:rotation-comparison}節の順序付き実表示と自然な関係直線同期を固定する。
$E_i/K$ は滑らかな genus-one 点曲線、$L_i$ は実次数三線束、
$p_i\in J_{E_i}(K)$ は点の平行移動であり、
\[
 Y=\prod_{i=1}^k E_i,\qquad
 R_{i,p}=T_{\OO_{E_i}}(t_{p_i})_*(-\otimes L_i),\qquad
 R_{i,0}=T_{\OO_{E_i}}(-\otimes L_i)
\]
と書く。関係直線は体上で一度自明化し、指定周期を
$\pi_{i,p}:R_{i,p}^3\simeq[2]$、
$\pi_{i,0}:R_{i,0}^3\simeq[2]$ と表示する。
この自明化を全因子で同時に変更しても以下の障害類は変わらない。

基準因子 $b$ を固定し、
\[
 G=Z_0=(\mu_3^k)/\Delta\mu_3,\qquad H=\widehat G,
 \qquad Q_{i,p}=R_{i,p}^{-1}R_{b,p},\quad
 Q_{i,0}=R_{i,0}^{-1}R_{b,0}\quad(i\ne b)
\]
とする。積上の関手には他の因子の恒等関手を外積するものとする。
$H$ の作用は指定周期を含む coherent な作用であり、
その同定を任意の $\operatorname{Aut}(H)$ で変更することはしない。
本節の全核という語は実際の perfect Fourier--Mukai 核を指す。
有限群の二重同変化により、これは元の残余圏の
$G$-同変 dg Morita 比較核の問題に等しい
\cite[Proposition~4.1, Theorem~4.2]{MEDElagin}。
体上の適用と rotation の同変性は
定理\paperref{main}{mbrc:thm:finite-point-comparison}の証明で確認済みである。

\subsection{外積分解と shift の曖昧さ}

\begin{lemma}[Orlov 像が block 対角である核の外積分解]
\label{mbfc:lem:factorization}
$\kappa$ を標数零の代数閉体とし、$E_1,\ldots,E_k$ を楕円曲線とする。
$\operatorname{FM}(E_i)$ と $\operatorname{FM}(Y)$ を
Fourier--Mukai 自己同値の同型類の群とし、
$\operatorname{FM}_{\mathrm{block}}(Y)$ を Orlov 像が
$Y\times\widehat Y=\prod_i(E_i\times\widehat E_i)$ に関して
block 対角である部分群とする。このとき外積は同型
\begin{equation}\label{mbfc:eq:factorization}
 \operatorname{FM}_{\mathrm{block}}(Y)
 \simeq
 \left(\prod_i\operatorname{FM}(E_i)\right)
 \big/
 \left\{([n_1],\ldots,[n_k]):\sum_i n_i=0\right\}.
\end{equation}
を与える。外積が恒等関手となる因子関手は、それぞれ shift である。
因子が互いに非同種なら $\operatorname{FM}_{\mathrm{block}}(Y)=\operatorname{FM}(Y)$ である。
\end{lemma}
\begin{proof}
Orlov の写像と完全列
\cite[Propositions~2.15, 3.3, Theorem~4.14]{MBFCOrlovAbelian}を用いる。
原点を選んで $Y$ をアーベル多様体とみなすと、
\[
 0\longrightarrow \mathbf Z\oplus(Y\times\widehat Y)(\kappa)
 \longrightarrow\operatorname{FM}(Y)
 \longrightarrow U(Y\times\widehat Y)\longrightarrow1
\]
がある。
isometric 自己同型の block 対角部分群は
\[
 U(Y\times\widehat Y)_{\mathrm{block}}=\prod_i U(E_i\times\widehat E_i)
\]
である。各成分を因子の Fourier--Mukai 同値へ持ち上げ、その外積を取る。
与えられた同値との差は上の完全列の核に入り、
平行移動・次数零線束の tensor・shift である。
平行移動と $\Pic^0(Y)=\prod_i\Pic^0(E_i)$ の線束も外積であるから全射性を得る。
非同種なら異なる因子とその双対の間の群準同型は零なので、全 Orlov 像が block 対角である。

核を確認する。因子同値の核 $P_i$ の外積が $\OO_{\Delta_Y}$ と同型なら、
Orlov の層性\cite[Proposition~3.2]{MBFCOrlovAbelian}により
$P_i=F_i[n_i]$ と書ける。
支持の外積が $\Delta_Y=\prod_i\Delta_{E_i}$ なので、各 $F_i$ の支持は
$\Delta_{E_i}$ である。対角のイデアルによる積も外積後に零となり、
非零な他因子との体上の tensor は忠実なので、各 $F_i$ は対角に scheme-theoretically 支持される。
外積が対角上の線束であることを、他の因子の点へ制限すると、
各 $F_i$ 自体が対角上の線束となる。
従って $P_i=\Delta_{E_i*}M_i[n_i]$ である。
同じ制限により各 $M_i$ は自明であり、cohomology の次数から
$\sum_i n_i=0$ を得る。逆の包含は外積の shift 則による。
\end{proof}

\subsection{指定された中心作用による三因子以上の分離}

ここでは代数閉体 $\kappa$、$\operatorname{char}\kappa=0$ 上で、
楕円曲線 $E_i$ と次数三線束 $L_i$ を考える。
因子間の同種性、CM の有無、または $j$ の相異なることは仮定しない。
原点 $o_i$ を選び、主偏極
\[
 \varpi_i:E_i\xrightarrow{\sim}\widehat E_i,\qquad
 x\longmapsto t_x^*\OO_{E_i}(o_i)\otimes\OO_{E_i}(o_i)^{-1}
\]
によって $V_i=E_i\times\widehat E_i$ を $E_i^2$ と表示する。
Orlov の写像を $\mathfrak u$ と書く。
その規約は
\[
 \Phi_{(x,\xi)}=t_{x*}(-)\otimes P_\xi,\qquad
 F\Phi_vF^{-1}\simeq\Phi_{\mathfrak u(F)v}
\]
である \cite[equation~(2.1), Corollary~2.13,
Propositions~2.15, 3.3]{MBFCOrlovAbelian}。

\begin{lemma}[因子 rotation の Orlov 行列]
\label{mbfc:lem:orlov-rotation}
上の $V_i=E_i^2$ の座標において、
$R_i=T_{\OO_{E_i}}\circ(-\otimes L_i)$ の Orlov 像は
\[
 r_i=r=
 \begin{pmatrix}-2&-1\\3&1\end{pmatrix},\qquad
 r^2+r+I=0.
\]
従って $r-I$ は isogeny であり、
$(r-I)(r+2I)=-3I$ である。
これは rank--degree 格子上の作用を代用したものではなく、
$E_i\times\widehat E_i$ 上の実際の群準同型の行列である。
\end{lemma}
\begin{proof}
まず $-\otimes L_i$ と $\Phi_{(x,\xi)}$ の共役を取る。
$\varpi_{L_i}=3\varpi_i$ なので、その作用は
$(x,\xi)\mapsto(x,\xi+3\varpi_i(x))$、従って行列は
\[
 t_3=\begin{pmatrix}1&0\\3&1\end{pmatrix}
\]
である。$\Pic^0$ 成分はこの共役行列に影響しない。
主偏極で元の曲線へ戻した normalized Poincare Fourier--Mukai 変換を
$\mathsf S$ とすると、skyscraper と次数零線束の像から
\[
 \mathfrak u(\mathsf S)=
   s=\begin{pmatrix}0&-1\\1&0\end{pmatrix}.
\]
実際、$\mathsf S(k_x)=P_{\varpi_i(x)}$、
$\mathsf S(P_{\varpi_i(y)})=k_{-y}[-1]$ である。
点の spherical twist は $T_{k_{o_i}}=-\otimes\OO(o_i)$、
$\mathsf S(k_{o_i})=\OO_{E_i}$ なので、
\[
 \mathfrak u(T_{\OO_{E_i}})
   =s\begin{pmatrix}1&0\\1&1\end{pmatrix}s^{-1}
   =\begin{pmatrix}1&-1\\0&1\end{pmatrix}.
\]
Orlov 写像の関手性から、この行列と $t_3$ の積が $r$ となる。
残りは整数行列の積で確認できる。
特に shift は Orlov 写像の核に入るので、
rank--degree 格子での shift の作用との同一視はしていない。
\end{proof}

$Y=\prod_{i=1}^kE_i$、$V=Y\times\widehat Y=\prod_iV_i$ とする。
基準 $b$ を選び、$Q_i=R_i^{-1}R_b$ とする。
$H=(\mathbf Z/3)^{k-1}$ の Orlov 像は
\[
 q_{\mathbf a}=\operatorname{diag}(r_1^{a_1},\ldots,r_k^{a_k}),
 \qquad
 \mathbf a\in(\mathbf Z/3)^k,\qquad \sum_i a_i=0
\]
である。

\begin{lemma}[三因子以上の整数 block rigidity]
\label{mbfc:lem:marked-block-rigidity}
$k\ge3$ とする。
$g\in U(Y\times\widehat Y)$ が全ての $q_{\mathbf a}$ と可換するなら、
$g$ は因子 $V_i$ に関して block 対角である。
従って
\[
 \operatorname{Cent}_{U(V)}(\mathfrak u(H))
   =\prod_i\operatorname{Cent}_{U(V_i)}(r_i).
\]
この主張は同種な因子や CM 係数を含む整数 Hom 群上で成立する。
\end{lemma}
\begin{proof}
$g_{ij}:V_j\to V_i$ を block 成分とする。
$i\ne j$ を固定し、これらと異なる $\ell$ を選ぶ。
$a_i=1$、$a_\ell=-1$、他の成分を零とした元について可換性を取れば
\[
 (r_i-I)g_{ij}=0.
\]
前補題の恒等式を左から掛けると $3g_{ij}=0$ となる。
アーベル多様体間の Hom 群は torsion-free なので $g_{ij}=0$ である。
具体的には、$3g_{ij}=0$ なら連結な $V_j$ の像は有限群 $V_i[3]$ に入り、
群準同型の像は零となる。
全ての非対角成分が零になった後、
isometry 条件と可逆性は因子ごとの条件となる。
逆包含は明らかである。
\end{proof}

補題\ref{mbfc:lem:factorization}により、ここで得た block 対角性は
実際の Fourier--Mukai 核の外積分解を与える。
同種因子間の混合線束は Neron--Severi 部分に属し、
Orlov 写像の核である $\Pic^0(Y)$ へ紛れ込むことはない。

ここで用いた Orlov の完全列と楕円曲線上の数値分類は既知の結果である。
本節で必要なのは、それらを指定された双対 rotation 作用へ適用する次の議論である。

\begin{proposition}[双対 rotation 作用の全 underlying centralizer]
\label{mbfc:prop:centralizer}
代数閉体上で、次のいずれかを仮定する：
\begin{equation}\label{mbfc:eq:block-hypothesis}
 k\ge3,\qquad\text{または}\qquad k=2\text{ かつ }E_1,E_2\text{ が非同種}.
\end{equation}
各 $L_i$ を次数三とし、$R_i=R_{i,0}$、
$\varepsilon_i=R_i^{-1}[1]$ と置く。
$\varepsilon_i^3=[1]$、$\varepsilon_i^2=R_i$ である。
全ての $Q_i=R_i^{-1}R_b$ と関手の同型類で可換する自己同値の群は
\begin{equation}\label{mbfc:eq:centralizer}
 \operatorname{Cent}_{\operatorname{FM}(Y)}\{Q_i:i\ne b\}
 \simeq
 \left(\prod_i\operatorname{Aut}(E_i,L_i)\right)\times\Lambda_k,
 \qquad
 \Lambda_k=\mathbf Z^k/
       \{(3n_i)_i:\sum_i n_i=0\}.
\end{equation}
ここで $(a_i)$ は外積 $\boxtimes_i\varepsilon_i^{a_i}$ を表す。
基準因子を選ぶと $\Lambda_k\simeq\mathbf Z\oplus(\mathbf Z/3)^{k-1}$ となる。
これは underlying 関手の同型類の分類であり、coherent な同変構造はまだ標識していない。
\end{proposition}
\begin{proof}
三ブロック以上では補題\ref{mbfc:lem:marked-block-rigidity}、
二つの非同種因子では補題\ref{mbfc:lem:factorization}により、
自己同値を $\boxtimes_i F_i$ と書く。
$FQ_iF^{-1}\simeq Q_i$ を右辺の逆関手と合成すると、
第 $i$、第 $b$ 因子だけを除いて恒等となる外積を得る。
同補題の核の記述から、ある整数 $n_i$ に対し
\[
 F_iR_iF_i^{-1}\simeq R_i[n_i]
\]
となり、基準因子にも同様の式が成り立つ。
三乗を取ると左辺は $[2]$、右辺は $[2+3n_i]$ である。
非零有界 coherent 複体は非零 shift と同型にならないため $n_i=0$ である。
従って全ての $F_i$ が $R_i$ と可換する。

系\ref{cor:nfm-elliptic-full-centralizer}により、各因子で
$F_i=\phi_{i*}\varepsilon_i^{a_i}$、
$\phi_i\in\operatorname{Aut}(E_i,L_i)$ と一意に書ける。
補題\ref{mbfc:lem:factorization}の shift の核は、
$\varepsilon_i^3=[1]$ の下で式\eqref{mbfc:eq:centralizer}の部分格子に一致する。
偏極自己同型は $R_i$ と shift に可換するので直積である。
\end{proof}

同じ議論を source の $R_{i,p}$ と target の $R_{i,0}$ に適用すると、
幾何的な underlying $H$-intertwiner は、各因子の全一因子 intertwiner の外積となる。
周期を三乗した上の議論は、因子間で余分な shift を相殺する可能性も除いている。

\subsection{双対群の抽象自己同型を許した正規化群}

この小節では $H$ の標識を $\operatorname{Aut}(H)$ によって変更することを許す。
まず関手の同型類の群における normalizer を扱い、
coherent な作用構造についての主張は明記する。
代数閉体 $\kappa$、$\operatorname{char}\kappa=0$ 上で
$Y=\prod_{i=1}^kE_i$、$R_i=T_{\OO_{E_i}}(-\otimes L_i)$ とし、
各 $L_i$ の次数は三、指定周期は $R_i^3\simeq[2]$ とする。
\[
 H=\left\{(t_i)\in\mathbf F_3^k:\sum_i t_i=0\right\}
\]
と表示する。整数 lift $\widetilde t_i$ に対し、$H$ の作用関手の同型類は
\begin{equation}\label{mbfc:eq:normalizer-action}
 Q_{\mathbf t}=
 \left(\mathop{\boxtimes}_{i=1}^k R_i^{\widetilde t_i}\right)
 \left[-\frac{2}{3}\sum_i\widetilde t_i\right]
\end{equation}
である。三次周期により lift の変更に依存しない。
特に $Q_{e_i-e_j}=R_iR_j^{-1}$ であり、
coherence は既存の $\Gamma$-同変 MF 模型からのものを用いる。
この作用は忠実である。

\begin{lemma}[三因子以上の指標配置]
\label{mbfc:lem:normalizer-characters}
$k\ge3$ とし、$\chi_i\in H^\vee$ を第 $i$ 座標の制限とする。
$\{\pm\chi_i:1\le i\le k\}$ を保つ $\operatorname{GL}(H^\vee)$ の部分群は
\[
 \{\pm1\}\times S_k
\]
である。すなわち、その任意の元は一つの置換と全因子共通の符号から成る。
$k=3$ にも例外はない。
\end{lemma}
\begin{proof}
$H^\vee=\mathbf F_3^k/\langle(1,\ldots,1)\rangle$ において
$\chi_i$ は標準基底の像である。
$k\ge3$ なら $\chi_i=\pm\chi_j$ は $i=j$ で符号が正の場合に限る。
実際、異なる $i,j$ に対して第三の座標を比較すれば、
$e_i\mp e_j$ が $(1,\ldots,1)$ の倍となることはない。
また、$\chi_i$ の間の関係全体は
\[
 \sum_i\chi_i=0
\]
の一次元空間である。
従って配置を保つ線型同型は $\chi_i\mapsto\epsilon_i\chi_{\sigma(i)}$
と書け、上の関係を移すと全ての $\epsilon_i\in\{\pm1\}$ が等しい。
逆に共通符号と任意の置換はこの関係を保ち、線型同型を定める。
これらの作用は忠実である。
特に $k=3$ では、四点からなる $\mathbf P^1(\mathbf F_3)$ のうち
選ばれた三点の集合を保つ射影作用は $S_3$ であり、
残る共通 scalar が上の $\{\pm1\}$ である。
\end{proof}

\begin{proposition}[全双対作用の Fourier--Mukai normalizer]
\label{mbfc:prop:normalizer}
$k\ge3$ とする。因子の非同種性、CM の有無、$j$ の相異なることは仮定しない。
Fourier--Mukai 自己同値の同型類の群において
\[
 N_H=\left\{F:\ \exists\alpha\in\operatorname{Aut}(H),\quad
 FQ_{\mathbf t}F^{-1}\simeq Q_{\alpha(\mathbf t)}\quad
 (\mathbf t\in H)\right\}
\]
と置く。
このとき $F\in N_H$ が誘導する $\alpha$ は因子置換だけであり、
共通反転は生じない。
許される置換は正確に
\[
 \Pi=\left\{\sigma\in S_k:
       (E_i,L_i)\simeq(E_{\sigma(i)},L_{\sigma(i)})\text{ for all }i\right\}
\]
であり、完全列
\begin{equation}\label{mbfc:eq:normalizer-sequence}
 1\longrightarrow\operatorname{Cent}_{\operatorname{FM}(Y)}(H)
 \longrightarrow N_H\longrightarrow\Pi\longrightarrow1
\end{equation}
を得る。
各 $F$ は置換後に因子 Fourier--Mukai 同値の外積となり、
その因子同値は $R_i$ を対応する $R_{\sigma(i)}$ へ shift なしで移す。

代数閉体上では、$\Pi$ は $j$ の値を保つ置換群でもある。
従って $j_i$ が相異なれば
\begin{equation}\label{mbfc:eq:normalizer-equals-centralizer}
 N_H=\operatorname{Cent}_{\operatorname{FM}(Y)}(H).
\end{equation}
同じ $j$ の因子ごとに偏極同型を選ぶと、
式\eqref{mbfc:eq:normalizer-sequence}は分裂し、
命題\ref{mbfc:prop:centralizer}の群に $\Pi$ が因子置換で作用する半直積となる。
この分裂は選択に依存する。
\end{proposition}
\begin{proof}
原点を選んで $V_i=E_i\times\widehat E_i$、$V=\prod_iV_i$ とする。
補題\ref{mbfc:lem:orlov-rotation}により、Orlov 写像での作用は
\[
 \mathfrak u(Q_{\mathbf t})
       =\operatorname{diag}(r^{t_1},\ldots,r^{t_k}),\qquad
 r=\begin{pmatrix}-2&-1\\3&1\end{pmatrix}.
\]
これは数値 Grothendieck 群の代用ではなく、実際の $V$ 上の作用である。
一つの素数 $\ell\ne3$ を固定し、Tate 加群を
$\mathbf Q_\ell(\zeta_3)$ へ拡大する。
$r$ の固有値は $\zeta_3,\zeta_3^{-1}$ なので、
$H$ の指標の集合は正確に $\{\pm\chi_i\}$ である。
各指標の重複度は二であり、因子の同種性や CM によって変わらない。

Orlov 写像の関手性
\cite[Proposition~2.15]{MBFCOrlovAbelian}から、
$u=\mathfrak u(F)$ は $\alpha$ に沿ってこの表現を移す。
従って前補題より、ある $\sigma\in S_k$ と共通符号
$\epsilon\in\{\pm1\}$ があり、
\[
 (\alpha\mathbf t)_{\sigma(i)}=\epsilon t_i
 \qquad(\mathbf t\in H)
\]
となる。

ここで有理的な指標分解から実際の整数 block 分解へ戻る。
各 $V_i$ の Tate 加群は共役指標対 $\{\chi_i,-\chi_i\}$ の成分である。
従って $u$ の $V_i\to V_j$ 成分は、$j\ne\sigma(i)$ なら
Tate 加群上で零となる。
アーベル多様体の準同型は Tate 加群への作用で忠実に検出されるので、
その成分は実際に零である。
これは Tate 同種定理を用いた全射性の主張ではなく、
非零な像を持つ準同型は Tate 加群上でも非零になるという忠実性だけを用いる。
同じ議論を $u^{-1}$ に適用すると、残る
\[
 u_i:V_i\xrightarrow{\sim}V_{\sigma(i)}
\]
は isogeny だけでなく実際の同型である。
Orlov の isometry 条件
\cite[Definition~2.17, Proposition~2.18]{MBFCOrlovAbelian}は
この block 分解に沿って各 $u_i$ の isometry 条件となる。

標数零の代数閉体上なので、各 $u_i$ は
$F_i:D^b(E_i)\xrightarrow{\sim}D^b(E_{\sigma(i)})$ へ持ち上がる。
具体的には、まず isometry の存在から一つの同値 $F_i^0$ を得て
\cite[Theorem~4.13]{MBFCOrlovAbelian}、その像を $u_i^0$ とする。
次に $(u_i^0)^{-1}u_i\in U(V_i)$ を同曲線の自己同値へ持ち上げ
\cite[Theorem~4.14]{MBFCOrlovAbelian}、$F_i^0$ と合成すれば、
指定された isometry $u_i$ 自身を持つ $F_i$ が得られる。
その外積と置換の合成を $F$ から除くと Orlov 写像の核に入り、
それは shift、平行移動、$\Pic^0(Y)$ の線束の tensor から成る
\cite[Proposition~3.3]{MBFCOrlovAbelian}。
$\Pic^0(Y)=\prod_i\Pic^0(E_i)$ なので、これらも因子の外積である。
従って $F$ 自身を、置換を伴う $F_i$ の外積としてよい。
ここでは異なる因子間の isometry の持上げを代数閉体上でだけ用いており、
任意の定義体上の全射性は仮定していない。

$\mathbf t=e_i-e_j$ の intertwining を取り、
外積が恒等になる核が総 shift 零だけであるという
補題\ref{mbfc:lem:factorization}の記述を適用する。
すると各因子について整数 $n_i$ があり、
\[
 F_iR_iF_i^{-1}\simeq R_{\sigma(i)}^{\epsilon}[n_i]
\]
となる。三乗を取れば
\[
 [2]\simeq[2\epsilon+3n_i]
\]
である。
非零有界 coherent 複体は非零 shift と同型にならない。
$\epsilon=-1$ なら $3n_i=4$ となって不可能であり、
$\epsilon=1$ なら $n_i=0$ である。
従って共通反転は生じず、各 $F_i$ は shift なしの rotation intertwiner となる。

系\paperref{main}{cor:nfm-elliptic-any-intertwiner-criterion}を $p=0$ に適用すると、
$F_i$ の存在から $(E_i,L_i)\simeq(E_{\sigma(i)},L_{\sigma(i)})$ が従う。
逆にこの偏極同型の pushforward の外積は所要の置換を実現する。
従って像は正確に $\Pi$ であり、核は定義から centralizer である。
代数閉体上では同じ $j$ の genus-one 曲線は同型であり、
次数三線束の差は平行移動で吸収できる。
実際、平行移動から $\Pic^3$ への写像の線型部は $[3]$ で、これは全射である。
同型類ごとの代表への偏極同型を固定すれば置換が厳密な群作用となり、
最後の分裂と半直積を得る。
\end{proof}

\begin{remark}[coherent 構造と符号の別検算]
\label{mbfc:rem:normalizer-coherence}
前命題は underlying 関手の同型類を分類する。
従って、coherent な $H$-normalizer に対しても誘導される自己同型の必要条件は同じである。
代数閉体上では、上で構成した各偏極置換の因子比較を三乗根で規格化して
指定周期に合わせ、その外積を取れば coherent な置換を実現できる。
一般の体上では period multiplier や比較核の降下を別に確認する必要がある。

共通反転の排除は CM の有無に依存しない関手の周期から得られた。
数値格子上でも、$Xr=r^{-1}X$ を満たす実行列は
\[
 X=\begin{pmatrix}a&b\\3b-3a&-a\end{pmatrix},\qquad
 \det X=-(a^2-3ab+3b^2)
\]
なので、非零なら行列式は負である。
従って $r$ と $r^{-1}$ は $\operatorname{SL}_2(\mathbf Z)$ では共役でない。
一方、上の証明では因子の Orlov isometry をこの二次元数値行列と同一視せず、
実際の Fourier--Mukai 持上げと周期を用いた。
\end{remark}

\begin{corollary}[二つの非同種因子と体上の位相下界]
\label{mbfc:cor:normalizer-bound}
$k=2$ で $E_1,E_2$ が幾何学的に非同種なら、
代数閉包上でも $N_H=\operatorname{Cent}_{\operatorname{FM}(Y)}(H)$ である。

さらに $K\supset\C$ 上の順序付き実表示を考え、次のいずれかを仮定する：
\begin{enumerate}
 \item $k\ge3$ かつ $j(E_i)$ が相異なる。
 \item $k=2$ かつ $E_1,E_2$ が幾何学的に非同種である。
\end{enumerate}
source の $H$ 作用を target の $H$ 作用へ
ある抽象自己同型 $\alpha\in\operatorname{Aut}(H)$ に沿って移す
coherent Fourier--Mukai 比較核があれば、必ず $\alpha=\mathrm{id}$ である。
従ってこれらの範囲では、中心 band の基底を忘れ、抽象自己同型を許しても、
指定された中心作用についての全核型判定は変わらない。
特に定理\ref{mbfc:thm:sharp-all-types}の独立な Hesse 族では、
この意味で標識を弱めても最小次数は $3^{k-1}$ と $3^k$ のままである。
\end{corollary}
\begin{proof}
二つの非同種因子の場合、任意の Fourier--Mukai 自己同値は外積となる。
$\operatorname{Aut}(H)=\{\pm1\}$ であり、反転を誘導する外積があれば、
前命題と同じく各因子で
$F_iR_iF_i^{-1}\simeq R_i^{-1}[n_i]$、従って $3n_i=4$ となる。
これで第一の主張を得る。

一般の $K$ 上では、まず代数閉包へ移す。
各 $p_i$ はそこで実 Picard 線束を持つので、既存の Picard 比較により
source の $H$ 作用を target の作用へ $\alpha=\mathrm{id}$ で移す基準核がある。
与えられた比較とその逆を合成すれば、target の $H$-normalizer になる。
前命題または第一の主張から $\alpha=\mathrm{id}$ である。
$H$ は定数群なので、この結論は元の $K$ 上でも同じである。
有限体拡大後にも仮定は保たれる。
従って全型判定とその Kummer 分裂次数に新しい核型は加わらない。
独立な Hesse 族の因子は幾何学的に非同種であることを
定理\ref{mbfc:thm:sharp-all-types}で既に確認している。
\end{proof}

この系は中心作用を比較データから完全に除く主張ではない。
normalizer に入らない残余圏自己同値の可能性は除いておらず、
完全に無標識な Kuznetsov 圏の全比較核や、その最小定義次数を分類していない。

\subsection{偏極自己同型が加える Kummer 位相}

\begin{proposition}[平面自己同型の period multiplier]
\label{mbfc:prop:multiplier}
$E/K$ を滑らかな平面三次、$L=\OO_E(1)$ とし、
$R_0=T_{\OO_E}(-\otimes L)$ とその三次方程式から来る指定周期
$\pi:R_0^3\simeq[2]$ を固定する。
$\phi\in\operatorname{Aut}_K(E,L)$ に対し、
$\eta:\phi_*R_0\simeq R_0\phi_*$ の指定周期への適合障害を
\[
 \nu_E(\phi)\in H^1(K,\mu_3)=K^\times/K^{\times3}
\]
と置く。これは群準同型である。

$W=H^0(E,L)$ とし、$\phi^*L\simeq L$ の選択から得る
pullback $g:W\xrightarrow{\sim}W$ を取る。
$f\in\operatorname{Sym}^3W$ を $E$ の方程式とし、
\begin{equation}\label{mbfc:eq:multiplier}
 g(f)=c_g f
\end{equation}
と書けば、pushforward の規約において
\begin{equation}\label{mbfc:eq:multiplier-class}
 \nu_E(\phi)=[c_g].
\end{equation}
$g$ の scalar 変更は $c_g$ を三乗倍するので右辺は選択に依存しない。
$R_0$ の冪、$[1]$ および $\varepsilon=R_0^{-1}[1]$ の multiplier は零である。
\end{proposition}
\begin{proof}
比較射の Hom は一次元であり、比較射を scalar 倍すると三重図式はその三乗倍になる。
二つの自己同型の比較射を合成すると三重図式の scalar は積となるため、
$\nu_E$ は群準同型である。

符号を確認する。$\phi_*$ に対応する graded 加群の制限は、
自由加群の行列に $g^{-1}$ を施す構成である。
従って MF の曲率は $f$ から $g^{-1}(f)=c_g^{-1}f$ へ移る。
この構成は内部 shift と MF の指定周期を保つ。
一方、命題\ref{prop:nfm-period-equation-frame-change}により
\[
 \pi_{c_g^{-1}f}=c_g\pi_f.
\]
従って同じ $\pi_f$ を source と target に使って三重図式を表示すると、
比較射の三乗の scalar は $c_g^{-1}$ である。
period-normalized 比較射の方程式は $v^3=c_g$ となるから、
式\eqref{mbfc:eq:multiplier-class}を得る。
この MF/window の比較が $\phi_*$ であることは、商の graded 加群上での制限と
section・evaluation による window の自然性から従う。

方程式を $u f$ に替えても両側の period が同じ $u^{-1}$ 倍を受けるので
$\nu_E$ は変わらない。
rotation の冪と shift については、MF の指標 shift と三次周期の中心 coherence を用いる。
それら自身の標準交換射は三乗周期を保つ。周期の cohomological degree は二なので、
$[1]$ との交換にも余分な符号は残らない。
\end{proof}

\begin{example}[基本楕円領域でも非零となる multiplier]
\label{mbfc:ex:nonzero-multiplier}
$K=\C(a)$ とし、$\mu\in\C$ を、Hesse 三次の判別式が非零かつ
$j\ne0,1728$ となるように選ぶ。
\[
 f=x^3+a y^3+a^2z^3-3\mu a xyz,
 \qquad g(x)=az,\quad g(y)=x,\quad g(z)=y
\]
と置くと $g^3=aI$、$g(f)=a f$ である。
代数閉包上で変数を三乗根により変更すれば、これは滑らかな Hesse 三次となる。
従って $g$ は $E=V(f)$ の実際の偏極自己同型を定め、
\[
 \nu_E(\phi)=[a]\ne0
\]
となる。非零性は $a$-進付値で分かる。
幾何的な $\operatorname{Aut}(E,0)$ は $\{\pm1\}$ であるから、
この位数三の自己同型は幾何的には非零三捩れ点の平行移動である。
一般の体上で全ての $\nu_E$ を零とみなすことはできない。
\end{example}

\subsection{全核型のスカラー障害と位相判定}

各因子の Picard gerbe と位相を第\ref{mbpd:sec:picard-descent}節の通り
\[
 \delta_i=\delta_{E_i}(p_i)\in H^2(K,\mathbb G_m),\qquad
 \gamma_i\in H^1(K,\mu_3),\qquad
 \alpha=\sum_i\delta_i
\]
と書く。$\gamma_i$ は、個々の $p_i$ が実線束で表されない場合にも
factor Picard gerbe 上の scalar 重みが零であることにより $K$ 上で定義される。
本節では $H^1(K,\mu_3)$ を加法記法で書く。

\begin{theorem}[因子分離の下での全中心同変核の存在判定]
\label{mbfc:thm:full-criterion}
上の順序付き実表示について、$k\ge3$、または
$k=2$ かつ $J_{E_1,\overline K},J_{E_2,\overline K}$ が非同種であるとする。
以下の条件は、$K$ の任意の体拡大後にも同じ形で成立する。

\begin{enumerate}
 \item 全ての underlying $H$-intertwiner の幾何的同型類は、
       $\boldsymbol\phi=(\phi_i)_i$、
       $\phi_i\in\operatorname{Aut}(E_{i,\overline K},L_i)$ と
       $\lambda\in\Lambda_k$ により一意に表示される。
       その型が $K$ 上に実際の核を持つためには $\phi_i$ が $K$ 上に定義される必要があり、
       このとき実現 gerbe の類は正確に $\alpha$ である。
 \item 中心 $G$ の作用を保つ通常の残余圏の比較核が $K$ 上に存在するための必要十分条件は
       \begin{equation}\label{mbfc:eq:ordinary-criterion}
        \alpha=0,\qquad
        \exists\,\phi_i\in\operatorname{Aut}_K(E_i,L_i):\quad
        \gamma_i+\nu_{E_i}(\phi_i)
           \text{ が全ての }i\text{ で等しい}
       \end{equation}
       である。
 \item さらに rotation 比較と指定三次周期を保つ核が存在するための必要十分条件は
       \begin{equation}\label{mbfc:eq:strong-criterion}
        \alpha=0,\qquad
        \exists\,\phi_i\in\operatorname{Aut}_K(E_i,L_i):\quad
        \gamma_i+\nu_{E_i}(\phi_i)=0\quad\text{for all }i
       \end{equation}
       である。
\end{enumerate}
ここで全ての shift と数値型を許している。
二つの同種因子の全算術的核型、中心作用そのものを忘れた核については主張していない。
\end{theorem}
\begin{proof}
まず代数閉包上で $[M_i]=p_i$ を満たす線束を選び、
$F_0=\boxtimes_i(-\otimes M_i)$ とする。
一因子比較により $F_0R_{i,p}\simeq R_{i,0}F_0$ がある。
従って任意の underlying $H$-intertwiner $F$ について、
$FF_0^{-1}$ は全ての $Q_{i,0}$ と可換する。
命題\ref{mbfc:prop:centralizer}により
\begin{equation}\label{mbfc:eq:normal-form}
 F\simeq C_{\boldsymbol\phi,\lambda}F_0,
 \qquad
 C_{\boldsymbol\phi,\lambda}
    =\boxtimes_i\bigl(\phi_{i*}\varepsilon_i^{a_i}\bigr),
 \qquad[(a_i)]=\lambda
\end{equation}
となり、$(\boldsymbol\phi,\lambda)$ は一意である。
$F_0$ の幾何的同型類は各 $p_i$ だけで決まり、Galois 不変である。
また $\varepsilon_i$ は $K$ 上の関手なので、
$F$ が $K$ 上の核なら一意性により各 $\phi_i$ は Galois 不変、従って $K$ 上へ降下する。
その偏極保存性も降下する。
実際、$\phi_i^*L_i\otimes L_i^{-1}$ は幾何的に自明な実線束であり、
固有基底変換により $K$ 上でも自明である。

従って $C_{\boldsymbol\phi,\lambda}$ は $K$ 上の実際の Fourier--Mukai 自己同値となる。
これの逆を $F$ と合成した核は、代数閉包上で対角上の線束である。
cohomology の集中、対角への支持、可逆層であることを忠実平坦降下すると、
その核は $\Delta_{Y*}M$、
\[
 [M]=\sum_i\operatorname{pr}_i^*p_i\quad\text{in }\Pic^0(Y)(K)
\]
となる。
逆にこの類の実線束があれば、式\eqref{mbfc:eq:normal-form}により全ての
$K$-有理な型の underlying 核が得られる。
命題\ref{mbpd:prop:product-gerbe}により、その実現 gerbe は
$\mathfrak G_{\mathrm{Pic}}$ と band を保って同値で、類は $\alpha$ である。
個々の $\delta_i$ の消滅は必要ではなく、ここではその和の相殺を保持している。

次に coherent な構造を調べる。
各 $\phi_i$ の period-normalized commutator の torsor を $V_i(\phi_i)$ と書く。
その類は $\nu_{E_i}(\phi_i)$ である。
$\varepsilon_i$ の標準比較は周期を保つので、
型 $(\boldsymbol\phi,\lambda)$ の各因子の位相 torsor は
\[
 T_i(\phi_i)=T_i\mathbin{\wedge^{\mu_3}}V_i(\phi_i),
 \qquad [T_i(\phi_i)]=\gamma_i+\nu_{E_i}(\phi_i)
\]
である。
これは因子線束が局所的にしか存在しない場合にも、Picard gerbe 上で計算して降下できる。
核の scalar 自己同型は比較射の両端で相殺し、追加の scalar 重みは生じない。

定理\ref{mbpd:thm:coherent-corrections}の証明をこの合成比較射に適用すると、
固定型に対する全 coherent $H$-構造の torsor は
\[
 D_{\boldsymbol\phi}
   =\left(\prod_i T_i(\phi_i)\right)/\Delta\mu_3,
 \qquad
 [D_{\boldsymbol\phi}]_b
   =\bigl(\gamma_b+\nu_{E_b}(\phi_b)
          -\gamma_i-\nu_{E_i}(\phi_i)\bigr)_{i\ne b}
\]
である。
これは構成した比較射だけの集合ではない。
一つの coherent 構造があれば、任意の別の構造との差は
$\operatorname{Hom}(H,\mathbb G_m)=G$ の指標となる。
実際、同値核の degree-zero endomorphism は
$\mathrm{HH}^0(Y)=H^0(Y,\OO_Y)=K$、負の自己 Ext は零であり、
生成元の比較の差は scalar、積の coherence はちょうど指標条件である。
従って上の torsor は全ての構造を尽くす。
rotation と指定周期まで保持すると $\prod_iT_i(\phi_i)$ が残る。
核の scalar 自己同型群はどちらの場合も $\mathbb G_m$ である。

以上の実現 gerbe と有限 torsor の両方が中立となる条件が、
それぞれ式\eqref{mbfc:eq:ordinary-criterion}、
式\eqref{mbfc:eq:strong-criterion}である。
二重同変化の自然性により元の $G$-同変残余比較へ移せる。
この議論は全ての体拡大上で同じなので、指定した任意の拡大体にも適用できる。
\end{proof}

\begin{remark}[全型という範囲]
\label{mbfc:rem:scope}
前定理は幾何的核の型、その各体拡大上の実現 gerbe、および coherent 比較の存在を扱う。
任意の非被約基底上の全核 stack を外積型へ分解する相対定理は主張していない。
また、例\ref{mbfc:ex:nonzero-multiplier}により、
Picard 補正型の $D$ が非自明であることだけから、
一般の $K$ 上で全ての中心同変比較核の不存在を推論することはできない。
全型では上の $\nu_{E_i}$ による補正を先に調べる必要がある。
\end{remark}

\subsection{Hesse 族における全型の最小次数}

\begin{lemma}[固定 Hesse level の全偏極自己同型]
\label{mbfc:lem:hesse-lifts}
$K\supset\C$ 上の滑らかな Hesse 三次
\[
 f=u(x^3+y^3+z^3)+vxyz,\qquad j(E)\ne0,1728
\]
を考える。
任意の体拡大 $K'/K$ 上の $\operatorname{Aut}_{K'}(E,L)$ の全ての元は、
$\zeta\in\C$ を原始三乗根として
\[
 A(x,y,z)=(x,\zeta y,\zeta^2z),\qquad
 B(x,y,z)=(y,z,x),\qquad
 S(x,y,z)=(y,x,z)
\]
の積による $\operatorname{GL}_3(\C)$ の lift を持つ。
それらは $f$ を厳密に保つので、全ての period multiplier は零である。
\end{lemma}
\begin{proof}
$o=(1:-1:0)$ を flex 原点とする。
この level の行列表示は
\cite[\S4, pp.~246--247]{NFMArtebaniDolgachev}の標準表示である。
$A,B$ は射影的に可換する位数三の元であり、相異なる九つの射影変換を生成する。
$j\ne0,1728$ より原点保存自己同型は $\{\pm1\}$ なので、
これらの位数三の元は三捩れ平行移動であり、$E[3]$ を尽くす。
$S$ は $o$ を固定する非恒等な involution なので inversion である。
次数三偏極を保存する曲線自己同型は
$E[3]\rtimes\{\pm1\}$ だから、この十八元が幾何的な群全体である。
全ての行列が $\C$ 上で定義されるため、任意の $K'$ 上でも追加の元はない。
$A,B,S$ は $x^3+y^3+z^3$ と $xyz$ をそれぞれ保つ。
命題\ref{mbfc:prop:multiplier}を適用すれば最後の主張を得る。
\end{proof}

\begin{theorem}[独立境界位相の全比較核に対する鋭い次数]
\label{mbfc:thm:sharp-all-types}
系\ref{mbpd:cor:phase-independence}の独立な $k$ 個の Hesse 族を取る。
すなわち Sklyanin パラメータを
\[
 (a_i,b_i,c_i)=(1+d_i,1,c_i),\qquad
 K_0=\C(d_1,\ldots,d_k)
\]
とし、$c_i\in\C$ は同系の滑らかさと基本楕円条件を満たす一般の定数とする。
source と target には自然な中心・点三次直線比較を使い、指定周期を固定する。
このとき
\begin{enumerate}
 \item 中心 $G$ の作用を保つ通常の残余比較核を得る有限体拡大の最小次数は
       $3^{k-1}$ である。
 \item さらに rotation 比較と指定三次周期を保つ核を得る有限体拡大の最小次数は
       $3^k$ である。
\end{enumerate}
ここでは Picard 補正型への制限を置かず、全ての dg Morita 比較核を許している。
同じ結論は $\operatorname{Frac}\C[[d_1,\ldots,d_k]]$ 上でも成り立つ。
中心作用を忘れた比較についての最小性は主張しない。
\end{theorem}
\begin{proof}
Hesse パラメータの三乗
\[
 \kappa_i=\frac{(a_i^3+b_i^3+c_i^3)^3}{a_i^3b_i^3c_i^3}
\]
は $a_i=1+d_i,b_i=1$ の下で $d_i$ の非定数有理関数である。
Hesse 三次の $j$ は $\kappa_i$ の非定数有理関数である
\cite[\S2, equations~(4)--(5)]{NFMArtebaniDolgachev}。
従って各因子の $j_i$ は $d_i$ だけの非定数有理関数である。
従って $j_1,\ldots,j_k$ は $\C$ 上代数的に独立である。
同種な二つの楕円曲線の $j$ は代数関係を満たすため、
因子は幾何学的に互いに非同種である。
後者の事実は、固定した $j$ の曲線の有限位数部分群とその商が
$\overline{\C(j)}$ 上に定義されることからも分かる。
完備局所体へ移しても、同種性があれば同じ代数関係が生じるので非同種性は保たれる。

Hesse flex により各 $p_i$ の実 Picard 線束があるので $\alpha=0$ である。
補題\ref{mbfc:lem:hesse-lifts}により、任意の有限体拡大上でも全ての
$\nu_{E_i}$ が零となる。
従って定理\ref{mbfc:thm:full-criterion}の全型判定は、
Picard 補正型の判定と一致する。

命題\ref{mbpd:prop:boundary-phase}と系\ref{mbpd:cor:phase-independence}により、
$K_0$ 上ではある $a_i(d_i)\in\C[d_i]_{(d_i)}^\times$ に対し
\[
 \gamma_i=[u_i],\qquad u_i=3d_i/a_i(d_i),\qquad
 v_{d_j}(u_i)=\delta_{ij}
\]
と書ける。
従って $\gamma_i$ の張る Kummer 部分空間の次元は $k$、
$\gamma_b-\gamma_i$ $(i\ne b)$ の張る部分空間の次元は $k-1$ である。
$K_0$ は $\mu_3$ を含むので、それぞれの全類を消す最小分裂体は
\[
 K_0\bigl(\sqrt[3]{u_b/u_i}:i\ne b\bigr),\qquad
 K_0\bigl(\sqrt[3]{u_i}:1\le i\le k\bigr)
\]
であり、その次数は $3^{k-1}$ と $3^k$ である。
全ての類を消す任意の有限体拡大はこれらの分裂体を含むので、この次数は下界でもある。
前定理がその拡大上の全型比較の存在を与えるため、最小性を得る。

完備局所環では全単元が三乗根を持つので $\gamma_i=[d_i]$ となる。
同じ付値計算と Kummer 理論を適用すれば、そこでの主張も従う。
ここで周期が境界まで延びる根拠は
系\paperref{main}{mbrn:cor:factor-boundary}の実 Koszul 最上 syzygy 直線を用いた一因子 window であり、
二捩れ境界で退化し得る $\det W_i$ への交代化同定は用いていない。
\end{proof}

\begin{corollary}[全中心同変核の大域切断の不存在]
\label{mbfc:cor:no-global-kernel}
前定理の有限型 Hesse 積族の基本楕円非空開上では、
中心 $G$ の作用を保つ通常の比較核の大域切断は存在しない。
任意の非空開へさらに縮小しても同じである。
\end{corollary}
\begin{proof}
大域核があれば一般点 $K_0$ へ引き戻せるが、
前定理の最小次数は $3^{k-1}>1$ である。
\end{proof}

\begin{remark}[既知の分類入力と本節の帰結]
\label{mbfc:rem:novelty}
アーベル多様体の Fourier--Mukai 完全列、楕円曲線の自己同値分類、
Hesse level と Kummer 理論は既知の入力である。
本節ではそれらを、満支持比較の指定された双対作用へ適用し、
因子 shift の相殺を除き、全核型の scalar 和と自己同型 multiplier による判定を得た。
Hesse 境界例では追加 multiplier が消えるため、
Picard 型で得た最小次数を全中心同変核へ拡張できる。
この適用・結論と同一の先行結果があるかどうかの新規性調査は未完了である。
\end{remark}

\subsection{二つの同種因子における混合中央化群}

以下では $k=2$、$Q=R_1^{-1}R_2$、
$q=\operatorname{diag}(r^{-1},r)$ とする。
主偏極を用いた Hom 行列の Rosati 転置を $X\mapsto X^\dagger$ とし、
\[
 J_2=\begin{pmatrix}0&1\\-1&0\end{pmatrix},\qquad
 X^\sharp=J_2^{-1}X^\dagger J_2
\]
と書く。全体では $J=\operatorname{diag}(J_2,J_2)$ を用いる。
Orlov の isometry 条件は $g^\sharp g=I$ である。
この式は $\operatorname{End}(E_i)=\mathbf Z$ を仮定していない。

\begin{proposition}[同種対に対する整数係数の混合元]
\label{mbfc:prop:two-isogenous-mixing}
$E_1,E_2$ が同種であれば、
$\operatorname{Cent}_{U(Y\times\widehat Y)}(q)$ は
非対角成分を持つ無限位数の元を含む。
この主張は $E_1\simeq E_2$ を必要とせず、CM の有無にも依存しない。
\end{proposition}
\begin{proof}
isogeny $u:E_2\to E_1$ を選び、その次数を $d$、
主偏極に関する dual isogeny を $u^\dagger$ とする。
$uu^\dagger=[d]$、$u^\dagger u=[d]$ である。
整数行列
\[
 A_0=\begin{pmatrix}-1&-1\\0&1\end{pmatrix},\qquad
 B_1=A_0,\qquad B_3=A_0(r-r^{-1})
     =\begin{pmatrix}-3&-1\\6&3\end{pmatrix}
\]
は、$s=1,3$ について
\[
 r^{-1}B_s=B_sr,\qquad
 rB_s=B_sr^{-1},\qquad
 B_s^\sharp=-B_s,\qquad B_s^2=sI
\]
を満たす。
$d$ と $3d$ の双方が平方数になることはないので、
$D=ds$ が非平方となる $s\in\{1,3\}$ を選べる。

$u,u^\dagger$ を二つの座標へ成分ごとに適用して
\[
 N=\begin{pmatrix}0&uB_s\\u^\dagger B_s&0\end{pmatrix}
\]
と置くと
\[
 Nq=qN,\qquad N^\sharp=-N,\qquad N^2=DI.
\]
Pell 方程式 $a^2-Db^2=1$ の $a>1,b>0$ となる整数解を取る。
この解の存在は非平方正整数に対する通常の Pell の定理による。
すると
\[
 g=aI+bN,\qquad g^{-1}=aI-bN
\]
は整数 Hom 成分を持つ実際のアーベル多様体自己同型であり、
$g^\sharp g=(a^2-Db^2)I=I$、$gq=qg$ を満たす。
二つの非対角成分は非零である。
$D$ は非平方なので $\mathbf Q[N]\simeq\mathbf Q(\sqrt D)$ であり、
この実二次体の埋込みの下で $g$ は
$a+b\sqrt D>1$ とその逆へ移る。従って $g$ は無限位数である。
\end{proof}

\begin{remark}[数値群の整数 Hom 表示]
\label{mbfc:rem:two-integral-centralizer}
二因子の数値中央化群は、$g=(g_{ij})$ に対して
\[
 rg_{ii}=g_{ii}r,\qquad
 r^{-1}g_{ij}=g_{ij}r\ (i\ne j),\qquad
 g^\sharp g=I
\]
を課した整数 Hom 行列の群である。
ここで $r^{-1}g_{ij}=g_{ij}r$ と $rg_{ij}=g_{ij}r^{-1}$ は同値である。
可換成分は $uI+vr$、反転 intertwiner は
$A_0(uI+vr)$ と一意に書ける。
係数 $u,v$ は各成分に対応する $\operatorname{Hom}(E_j,E_i)$ に属する。
これは $r$ の $(1,2)$ 成分が $-1$ であることを使って
四つの成分方程式を解けば分かる。
従って有理化で失われる整除条件を隠しておらず、
CM を含む Hom 群のままでの記述となっている。
\end{remark}

\begin{lemma}[数値中央化元の実際の coherent 核への持上げ]
\label{mbfc:lem:cyclic-centralizer-lift}
代数閉体 $\kappa$ 上で、任意の
$g\in\operatorname{Cent}_{U(Y\times\widehat Y)}(q)$ は
指定された $\langle Q\rangle\simeq\mathbf Z/3$ 作用と
coherently 可換する実際の Fourier--Mukai 自己同値へ持ち上がる。
特に前命題の混合元は、単なる数値的候補ではない。

underlying 自己同値の同型類については正確列
\[
 1\longrightarrow
   \mathbf Z\times E_1[3](\kappa)\times E_2[3](\kappa)
 \longrightarrow \operatorname{Cent}_{\operatorname{FM}(Y)}(Q)
 \xrightarrow{\mathfrak u}
   \operatorname{Cent}_{U(Y\times\widehat Y)}(q)
 \longrightarrow1
\]
がある。固定した underlying 中央化関手上の coherent 構造の同型類は
$\mu_3$-torsor である。
\end{lemma}
\begin{proof}
Orlov の全射 \cite[Theorem~4.14]{MBFCOrlovAbelian}で $g$ を $F$ へ持ち上げる。
可換子の Orlov 像は恒等なので、ある $v\in Y\times\widehat Y$ と整数 $n$ に対し
\[
 FQF^{-1}\simeq\Phi_v[n]\,Q
\]
と書ける。$Q^3\simeq\mathrm{id}$ の指定された関係を三乗に用いると、
$I+q+q^2=0$ により右辺の三乗は $[3n]$ となる。
有界 coherent 圏上の非零 shift は恒等関手ではないので $n=0$ である。

$(1-q)(2+q)=3I$ から $1-q$ は isogeny であり、
代数閉体上で全射である。
$(1-q)w=-v$ となる $w$ を取り、$F'=\Phi_wF$ と置けば
\[
 F'Q(F')^{-1}
   \simeq \Phi_{w+v-qw}\,Q
   \simeq Q.
\]
従って underlying 中央化関手が得られる。
この補正の符号は、$\Phi_w$ を $F$ の左から合成する規約による。

同型 $\eta:F'Q\xrightarrow{\sim}QF'$ を一つ選ぶ。
$Q^3\simeq\mathrm{id}$ の指定された coherence を使って
$\eta$ を三回合成した図式は、$F'$ の scalar 自己同型 $c$ となる。
実際、可逆核の零次自己準同型は
$\mathrm{HH}^0(Y)=H^0(Y,\OO_Y)=\kappa$ であり、負の自己 Ext は零である。
$\eta$ を $t\eta$ に替えるとこの scalar は $t^3c$ になる。
$t^3=c^{-1}$ を選べば、巡回群の三乗関係を保つ比較を得る。
負の自己 Ext の消滅により比較の mapping space に更なる高次自己同型はなく、
指定された群作用の coherence とこの三乗関係が実際の coherent intertwiner を与える。
二つの選択の差は
$\operatorname{Hom}(\mathbf Z/3,\kappa^\times)=\mu_3$ である。

最後に Orlov 写像の核の $\Phi_w[n]$ が $Q$ と可換する条件は $qw=w$ である。
$r$ と $r^{-1}$ の固定部分は、上の主偏極座標でともに
$\{(x,0):x\in E_i[3]\}$ である。
これと shift の独立性から正確列を得る。
\end{proof}

ここまでの二因子の持上げは代数閉体上の主張である。
一般の定義体上で数値中央化元の全てが実核を持つことや、
その各混合型の scalar gerbe・周期位相を分類したとはしていない。
一方、同種な二因子について外積型だけで全核を尽くすことはできない、
という障害は実際の coherent 中央化核によって証明されている。

整数行列と Pell 元の有限検算は
\texttt{computations/verify\_multiblock\_isogenous.py} に収める。
このコードは Orlov の存在定理、忠実平坦降下、または coherent 核の持上げを
機械的に証明するものではない。

\begin{example}[二次同種で生じる混合反転核]
\label{mbfc:ex:two-isogenous-inversion}
二因子の場合、上の非同種性を単に $j_1\ne j_2$ へ弱めることはできない。
代数閉体 $\kappa$、$\operatorname{char}\kappa=0$ 上で
$\phi:E_1\to E_2$ を次数二の isogeny とし、主偏極による dual を $\phi^\vee$ とする。
各因子の任意の次数三線束から $R_i$ を作り、
$Q=R_1^{-1}R_2$、$D=\operatorname{diag}(r^{-1},r)$ と置く。
整数行列
\[
 s=\begin{pmatrix}1&0\\-3&-1\end{pmatrix},\qquad
 J_2=\begin{pmatrix}0&1\\-1&0\end{pmatrix}
\]
は
$s^2=I$、$sr=r^{-1}s$、$s^{\mathsf t}J_2s=-J_2$ を満たす。
$\phi,\phi^\vee$ を二つの座標へ成分ごとに作用させ、
\begin{equation}\label{mbfc:eq:two-inversion-matrix}
 U=\begin{pmatrix}s&\phi^\vee I_2\\\phi I_2&s\end{pmatrix}
   :E_1^2\times E_2^2\longrightarrow E_1^2\times E_2^2
\end{equation}
と置く。このとき
\[
 U^{-1}=\begin{pmatrix}-s&\phi^\vee I_2\\\phi I_2&-s\end{pmatrix},
 \qquad U^\dagger\operatorname{diag}(J_2,J_2)U
          =\operatorname{diag}(J_2,J_2),\qquad UD=D^{-1}U.
\]
ここで $\dagger$ は主偏極に関する Rosati 転置を含む転置である。
これらは $\phi^\vee\phi=[2]$、$\phi\phi^\vee=[2]$ と
$s^{\mathsf t}J_2+J_2s=0$ から直接従う。
従って $U$ は整数 Hom 成分を持つ実際の Orlov isometry であり、
二つの非対角成分は非零である。

これは実際の coherent 反転核へ持ち上がる。
Orlov の定理\cite[Theorem~4.14]{MBFCOrlovAbelian}で $U$ を $F$ へ持ち上げると、
ある $z\in Y\times\widehat Y$、整数 $n$ に対し
\[
 FQF^{-1}\simeq\Phi_z Q^{-1}[n]
\]
と書ける。$I+D+D^2=0$ と $Q^3\simeq\mathrm{id}$ を用いて三乗すると
$[3n]\simeq\mathrm{id}$、従って $n=0$ である。
$1-D^{-1}$ は isogeny なので、代数閉体上で
$(1-D^{-1})w=-z$ となる $w$ が存在する。
$F'=\Phi_wF$ と置けば
\[
 F'Q(F')^{-1}\simeq\Phi_{z+w-D^{-1}w}Q^{-1}\simeq Q^{-1}.
\]
同型 $F'Q\simeq Q^{-1}F'$ の三乗図式のずれは scalar であり、
その逆数の三乗根で比較射を規格化すれば
$\mathbf Z/3$ の指定された関係を保つ。
可逆核の自己準同型は $\kappa$、負の自己 Ext は零なので、
これは実際の coherent $H$-normalizer を与える。
この補正は Orlov 像を変えないため、得られた核は因子の外積ではない。

一般の二次同種対では $j(E_1)\ne j(E_2)$ となるので、
異なる $j$ だけでは二因子の反転を排除できない。
この例では代数閉体上の存在を述べている。
一般の定義体上への核・coherence の降下や、
個々の因子 rotation・指定周期まで保つ比較、
混合型の scalar gerbe と Kummer 位相は結論していない。
\end{example}
